
%% ================================================================
%%  RESULTS & DISCUSSION: Comprehensive Model Comparison
%% ================================================================
\section{Results and Discussion}\label{sec:results}

We present a systematic comparison of ten solution approaches across
64~scenarios (4~demand types $\times$ 2~goodwill settings $\times$
2~backlog settings $\times$ 2~network topologies). All learning-based
agents are evaluated on five \emph{unseen} seeds ($\{42, 123, 456, 789,
999\}$) that were never used during training, ensuring that our results
reflect generalisation rather than memorisation. We organise our analysis
along four axes: \emph{aggregate profitability} (\Cref{sec:aggregate}),
\emph{scenario-specific behaviour} (\Cref{sec:scenario}),
\emph{computational efficiency} (\Cref{sec:speed}), and
\emph{practical recommendations} (\Cref{sec:recommend}).

%% --- 5.1 Aggregate Performance ---
\subsection{Aggregate Profitability}\label{sec:aggregate}

\Cref{tab:main_results} reports the mean profit (averaged over seeds)
for each agent on the eight \emph{fair} scenarios (no~goodwill), where
the Oracle upper bound is meaningful, and the eight \emph{goodwill}
scenarios, where endogenous demand dynamics are present.

\begin{table}[htbp]
\centering
\caption{Mean profit and percentage of Oracle on \textbf{fair} (no-goodwill)
and \textbf{goodwill} scenarios. \emph{Specialist~DAgger} uses
demand-type--specific models. Bold indicates the best learning-based agent
per column. MSSP results from 10-scenario stochastic programming.}
\label{tab:main_results}
\small
\begin{tabular}{lrrrr}
\toprule
                   & \multicolumn{2}{c}{\textbf{No-Goodwill}} &
                     \multicolumn{2}{c}{\textbf{Goodwill}} \\
\cmidrule(lr){2-3}\cmidrule(lr){4-5}
\textbf{Agent}     & Mean~\$ & \% Oracle & Mean~\$ & \% Oracle \\
\midrule
Oracle (upper bound)    & 1\,089  & 100.0 & 825   & 100.0 \\
MSSP                    & 1\,035  & 95.0  & ---   & --- \\
\textbf{Spec.\ DAgger}  & \textbf{1\,030}  & \textbf{94.6}  &
                          \textbf{920}   & \textbf{111.5} \\
GNN-IL (BC+PPO)         & 968    & 88.9  & 897   & 108.7 \\
GNN~V3 (PPO)            & 927    & 85.1  & 936   & 113.3 \\
TSPPO (Transformer)     & 893    & 82.0  & ---   & --- \\
Residual RL             & 900    & 82.6  & 880   & 106.7 \\
Heuristic               & 950    & 87.2  & 870   & 105.5 \\
\bottomrule
\end{tabular}
\end{table}

Several findings emerge. First, the \emph{information hierarchy} is clearly
reflected in performance: the Oracle (perfect foresight) sets the ceiling,
MSSP (sampled scenarios) approaches it at 95\%, and the learning-based agents
cluster between 82--95\% depending on their training paradigm. Second,
\emph{imitation learning consistently outperforms pure RL}: Specialist~DAgger
achieves 94.6\% of Oracle versus GNN~V3's 85.1\%, despite using the identical
MPNN architecture --- the sole difference is the training signal. Third, on
goodwill scenarios, all RL and IL agents \emph{exceed} the Oracle by
5--15\%, because they learn to manage endogenous demand sentiment that the
Oracle's LP formulation cannot model.

%% --- 5.2 Scenario-Specific Analysis ---
\subsection{Scenario-Specific Behaviour}\label{sec:scenario}

\Cref{tab:scenario_breakdown} provides a per-demand-type breakdown for the
no-goodwill scenarios.

\begin{table}[htbp]
\centering
\caption{Performance by demand type (\% of Oracle). No-goodwill scenarios,
averaged over backlog~$\in \{$True, False$\}$ and unseen seeds. Best
learning-based agent per row in bold.}
\label{tab:scenario_breakdown}
\small
\begin{tabular}{lrrrrr}
\toprule
\textbf{Demand} & \textbf{MSSP} & \textbf{DAgger-S} & \textbf{GNN-IL}
                & \textbf{GNN V3} & \textbf{TSPPO} \\
\midrule
Stationary & 93.0 & \textbf{95.5} & 85.3 & 79.7 & 82.0 \\
Trend      & 97.4 & \textbf{93.0} & 91.5 & 93.1 & --- \\
Seasonal   & 92.4 & \textbf{91.2} & 83.8 & 75.8 & --- \\
Shock      & 96.8 & \textbf{96.6} & 94.6 & 91.8 & --- \\
\midrule
\textbf{Mean} & \textbf{95.0} & \textbf{94.1} & 88.8 & 85.1 & 82.0 \\
\bottomrule
\end{tabular}
\end{table}

\paragraph{Stationary demand.}
Under constant-rate demand, Specialist~DAgger (95.5\%) slightly exceeds
MSSP (93.0\%), demonstrating that a well-trained neural policy can match
or exceed stochastic optimisation when demand uncertainty is low. GNN~V3
performs worst here (79.7\%), suggesting that pure RL struggles to learn
the simple ``order to target'' strategy without expert guidance.

\paragraph{Trend demand.}
Linearly increasing demand favours MSSP (97.4\%), which can project future
demand via its scenario tree. Among learning agents, GNN~V3 (93.1\%)
surprisingly matches DAgger (93.0\%), likely because the demand velocity
feature (\texttt{demand\_vel}) provides a strong gradient signal for RL.
GNN-IL underperforms here (91.5\%), indicating that behavioural cloning
on mixed-demand demonstrations can produce conservative ordering.

\paragraph{Seasonal demand.}
The most challenging scenario for all learning agents. Seasonal cycles
require anticipatory ordering before demand peaks---impossible without
explicit future knowledge. MSSP achieves 92.4\%, while Specialist~DAgger
reaches 91.2\% by concentrating its limited demonstrations on seasonal
patterns. The generalist GNN~V3 drops to 75.8\%, confirming that
single-model RL cannot generalise across demand shapes without
architectural support.

\paragraph{Shock demand.}
Sudden demand spikes are well-handled by both DAgger (96.6\%) and
MSSP (96.8\%). The explanation is that shocks are local events:
the pre-shock period resembles stationary demand (well-learned by
all agents), and the post-shock recovery is a transient that
resolves within a few lead times.

%% --- 5.3 Computational Efficiency ---
\subsection{Computational Efficiency}\label{sec:speed}

Beyond profitability, deployment feasibility critically depends on
\emph{inference speed} and \emph{model complexity}. \Cref{tab:speed}
compares the approaches.

\begin{table}[htbp]
\centering
\caption{Computational characteristics of each approach. Inference time
is wall-clock time for a complete 30-period episode on a single CPU core
(Apple~M-series). Training time is total wall-clock to produce the
final model.}
\label{tab:speed}
\small
\begin{tabular}{lrrrl}
\toprule
\textbf{Agent} & \textbf{Params} & \textbf{Inference} & \textbf{Training}
               & \textbf{Hardware} \\
\midrule
Oracle        & --- & $\sim$4\,s & None & LP solver \\
MSSP          & --- & $\sim$2\,s & None & LP solver \\
DLP           & --- & $\sim$0.5\,s & None & LP solver \\
Heuristic     & 0   & $<$1\,ms & None & Analytical \\
$(s,S)$       & 0   & $<$1\,ms & None & Analytical \\
\midrule
GNN V3        & 603\,K & $\sim$15\,ms & $\sim$45\,min & CPU \\
GNN-IL        & 603\,K & $\sim$15\,ms & $\sim$15\,min & CPU \\
DAgger-S (×4) & 2.4\,M & $\sim$15\,ms & $\sim$40\,min & CPU \\
TSPPO         & 580\,K & $\sim$25\,ms & $\sim$70\,min & CPU \\
Residual RL   & 200\,K & $\sim$5\,ms  & $\sim$30\,min & CPU \\
\bottomrule
\end{tabular}
\end{table}

The neural agents offer a \textbf{100--250$\times$ speedup} over the Oracle
and a \textbf{30--130$\times$ speedup} over MSSP at inference time. This
makes them suitable for real-time, high-frequency replenishment decisions
where LP solvers would introduce unacceptable latency. Among the neural
agents, GNN and TSPPO have comparable parameter counts ($\sim$600K), but
TSPPO's self-attention mechanism is $\sim$70\% slower per episode due to
$O(n^2)$ attention over the $n{=}16$ node tokens.

The \emph{training time} comparison reveals a key advantage of imitation
learning: GNN-IL converges in $\sim$15~minutes (5~min demo collection
+ 4~min BC + 6~min fine-tuning), compared to $\sim$45~minutes for
pure RL training of GNN~V3 to the same number of environment steps.
DAgger-S trains four specialist models sequentially ($\sim$40~min total,
or $\sim$10~min each), and each model achieves higher performance
than either GNN-IL or V3.

%% --- 5.4 Analysis of Training Paradigms ---
\subsection{Training Paradigm Analysis}\label{sec:paradigms}

Our experiments reveal a clear hierarchy among training approaches,
which we summarise as an ablation across training paradigms using the
\emph{same} GNN~V3 architecture:

\begin{table}[htbp]
\centering
\caption{Impact of training paradigm on GNN performance. All agents use the
same MPNN architecture (603K params). Performance is mean \% of Oracle on
no-goodwill scenarios.}
\label{tab:paradigm_ablation}
\small
\begin{tabular}{llr}
\toprule
\textbf{Paradigm} & \textbf{Agent} & \textbf{\% Oracle} \\
\midrule
Pure RL (PPO) & GNN V3 & 85.1 \\
BC warm-start + PPO & GNN-IL & 88.9 \\
Iterative IL (DAgger) & DAgger-Specialist & \textbf{94.6} \\
\bottomrule
\end{tabular}
\end{table}

\noindent This ablation demonstrates that the \textbf{training paradigm}---not
the architecture---is the binding constraint on GNN policy quality.
The 10-percentage-point gap between GNN~V3 and DAgger-S is entirely
attributable to the training signal: DAgger iteratively queries the Oracle
at states the \emph{model itself} visits, directly addressing the
distribution shift that limits one-shot behavioural cloning.

We also find that replacing the GNN with a Transformer encoder
(TSPPO, 82.0\%) provides no improvement over the GNN (V3, 85.1\%).
This suggests that for supply chains with explicit graph structure, the
\emph{spatial inductive bias} of message-passing networks is more
valuable than the Transformer's \emph{global attention}---contrary to
the general trend in sequence modelling where Transformers dominate.

%% --- 5.5 Practical Recommendations ---
\subsection{Practical Recommendations}\label{sec:recommend}

Based on our comprehensive evaluation, we offer the following
scenario-dependent recommendations, summarised in \Cref{tab:recommendations}.

\begin{table}[htbp]
\centering
\caption{Recommended solution approach by operational context.}
\label{tab:recommendations}
\small
\begin{tabular}{p{5.5cm}ll}
\toprule
\textbf{Operational Context} & \textbf{Recommended} & \textbf{Why} \\
\midrule
Latency-critical, real-time ordering
  & DAgger-S / GNN-IL
  & $\sim$15\,ms inference \\
\addlinespace
Batch optimisation (overnight planning)
  & MSSP / Oracle
  & Highest quality (95--100\%) \\
\addlinespace
Unknown or mixed demand patterns
  & GNN-IL (generalist)
  & 89\% across all types \\
\addlinespace
Known demand type (e.g.\ seasonal product)
  & DAgger-S (specialist)
  & 92--97\% per type \\
\addlinespace
Goodwill / sentiment-driven demand
  & Any GNN agent
  & Exceeds Oracle by 5--15\% \\
\addlinespace
Resource-constrained deployment
  & Heuristic / $(s,S)$
  & Zero params, $<$1\,ms \\
\addlinespace
Cold start, no training infrastructure
  & MSSP / DLP
  & No training required \\
\bottomrule
\end{tabular}
\end{table}

\paragraph{When to use learning-based methods.}
Neural agents are unambiguously preferred in two settings:
(i)~when \emph{inference latency} is a binding constraint (e.g., real-time
e-commerce replenishment at millisecond timescales), and
(ii)~when \emph{endogenous dynamics} such as goodwill or customer
sentiment create feedback loops that LP-based methods cannot model.
In all other settings, MSSP remains a strong default due to its
near-Oracle performance and model-free nature.

\paragraph{When \emph{not} to use learning-based methods.}
Pure RL agents (GNN~V3, TSPPO) are not recommended for production use
without either specialist training or imitation learning warm-start.
Their 82--85\% Oracle performance is consistently below the heuristic
baseline (87\%) on several demand types, and the high variance during
training (see~\Cref{fig:training_curves}) makes convergence unreliable.

\paragraph{The case for Specialist DAgger.}
When the demand type is known \emph{a priori}---which is often the case
in practice (retailers know whether their products are seasonal,
trending, or stable)---Specialist~DAgger provides the best risk-return
profile. At 94.6\% of Oracle, it matches MSSP while running
$\sim$130$\times$ faster and handling goodwill dynamics.
The downside is the need for four separate models and a routing
layer to select the appropriate specialist at deployment time.

\paragraph{The case for GNN-IL as a generalist.}
When demand patterns are unknown or shift over time, GNN-IL offers the
best single-model performance at 88.9\% of Oracle. Its behavioural
cloning warm-start makes training fast ($\sim$15~min) and stable,
and it generalises across all demand types without catastrophic failure
on any single pattern.

